% !TeX root = ../complete.tex
\section{مدل‌های تأخیر برای \lr{Timing Analysis}}
\label{sec:delay-timing-models}

برای مقایسه اولیه توپولوژی‌ها، مدل \lr{Logical Effort}
با چند پارامتر، اثر بار و نوع گیت را نشان می‌دهد. در تحلیل
زمان‌بندی یک شبکه، گیت‌ها ورودی‌هایی با زمان گذار متفاوت
دریافت می‌کنند و ظرفیت بار و جهت گذار نیز از یک گره تا گره
بعدی تغییر می‌کند. تحلیل‌گر باید علاوه بر تأخیر، زمان گذار
خروجی را تعیین کند، زیرا همان شکل موج ورودی مرحله بعد است
\cite[بخش \secnum{۴.۶}، صفحات ۱۷۳ و ۱۷۴]{weste}.

\subsection{گذار ورودی و مدل خطی وابسته به آن}
\label{sec:delay-slew-linear}

در مدل تأخیر یک \lr{Standard Cell} ترکیبی، مسیر زمانی از یک پایه ورودی
تا یک پایه خروجی، \lr{Timing Arc} نام دارد. برای مثال، در
\lr{Inverter}، صعود ورودی موجب افت خروجی می‌شود. زمان گذار
ورودی را با \(s_{\mathrm{in}}\) نشان می‌دهیم. این کمیت
مدت عبور شکل موج ورودی بین دو سطح مرجع است و در اینجا واحد
زمان دارد. تعریف و مرزهای اندازه‌گیری آن را باید همراه داده
کتابخانه خواند. زمان گذار از
\hyperref[sec:delay-waveform]{تأخیر میان عبور ورودی و خروجی}
متمایز است.

بخش \secnum{۴.۶.۱} کتاب، مدل خطی را با جمله‌ای برای اثر
زمان گذار ورودی گسترش می‌دهد. در صورت‌بندی کتاب، تأخیر مرحله
قبلی نماینده زمان گذار ورودی است. اگر \(t_{\mathrm{prev}}\)
تأخیر زمانی مرحله قبلی و \(q\) جهت گذار خروجی باشد،
مدل کتاب را می‌توان چنین نوشت
\cite[بخش \secnum{۴.۶.۱}، صفحه ۱۷۳]{weste}.
\begin{equation}
  t_{\mathrm{pd},q}=t_{0,q}+R_q C_L+k_qt_{\mathrm{prev}}.
  \label{eq:delay-slew-linear}
\end{equation}
پارامتر \(t_{0,q}\) تأخیر ذاتی در مدل و \(R_q\) ضریب اثر
بار خازنی \(C_L\) است. \(R_q\) واحد مقاومت دارد، زیرا
حاصل‌ضرب آن در ظرفیت باید واحد زمان داشته باشد. \(k_q\) بی‌بعد
است. برای صعود و افت خروجی، مجموعه ضرایب جداگانه به کار
می‌رود. رابطه \eqref{eq:delay-slew-linear} در بازه‌ای محدود
معتبر است، زیرا
تناسب زمان گذار ورودی با تأخیر مرحله قبلی در همه شکل موج‌ها
یکسان نمی‌ماند و یک خط ممکن است دامنه بزرگی از بارها و
زمان‌های گذار را به‌خوبی توصیف نکند.

\subsection{مدل جدولی و درون‌یابی تأخیر}
\label{sec:delay-nldm}

در \lr{Nonlinear Delay Model (NLDM)}، تأخیر از روی
جدولی با دو ورودی \(s_{\mathrm{in}}\) و \(C_L\) به دست
می‌آید. برای هر \lr{Timing Arc}، تأخیر و زمان گذار
خروجی در جهت‌های صعود و افت از جدول‌های جدا خوانده می‌شوند.
داده‌های زمانی \lr{Standard Cell}ها در قالب \lr{Liberty}
و معمولاً در فایل‌های \code{.lib} عرضه می‌شوند. داده باید
متناسب با شرایط مشخص فناوری، تغذیه و دما انتخاب شود
\cite[بخش‌های \secnum{۴.۶} و \secnum{۴.۶.۲}، صفحات
۱۷۳ و ۱۷۴]{weste}.

جدول \ref{tab:delay-nldm-interpolation} بخشی از داده‌های
افت خروجی \lr{Inverter} را نشان می‌دهد. می‌خواهیم تأخیر
را برای زمان صعود ورودی \qty{30}{\pico\second} و بار
\qty{3}{\femto\farad} تخمین بزنیم. نقطه موردنظر در
میانه هر دو بازه است، پس برای درون‌یابی خطی در هر راستا،
دو سر بازه وزن برابر دارند.

\begin{table}[htbp]
  \centering
  \caption{تأخیر افت خروجی \(t_{\mathrm{pdf}}\) برحسب
  \unit{\pico\second}، برای دو بار \(C_L\) و دو زمان صعود
  ورودی \(s_{\mathrm{in}}\).
  داده‌ها از جدول \secnum{۴.۷} کتاب، صفحه ۱۷۴، انتخاب شده‌اند
  \cite{weste}.}
  \label{tab:delay-nldm-interpolation}
  \begin{tabular}{ccc}
    \toprule
    \(C_L\) برحسب \unit{\femto\farad}
      & \(s_{\mathrm{in}}=\qty{20}{\pico\second}\)
      & \(s_{\mathrm{in}}=\qty{40}{\pico\second}\) \\
    \midrule
    \num{2} & \num{20.2} & \num{24.1} \\
    \num{4} & \num{33.8} & \num{37.6} \\
    \bottomrule
  \end{tabular}
\end{table}

ابتدا در هر ردیف بین دو زمان صعود ورودی میانگین می‌گیریم.
\(t_2\) و \(t_4\) تخمین‌های تأخیر در بار
\qty{2}{\femto\farad} و \qty{4}{\femto\farad} هستند.
سپس میان این دو نتیجه در راستای بار میانگین می‌گیریم.
\[
  \begin{aligned}
    t_2&=\frac{20.2+24.1}{2}\,\unit{\pico\second},&
    t_4&=\frac{33.8+37.6}{2}\,\unit{\pico\second},\\
    t_{\mathrm{pdf}}
      &\approx\frac{t_2+t_4}{2}
       \approx\qty{28.9}{\pico\second}.
  \end{aligned}
\]
این محاسبه نمونه‌ای از درون‌یابی دوبعدی است. در حالت کلی،
وزن‌ها از فاصله نقطه موردنظر تا دو سر هر بازه تعیین می‌شوند.
عدد به‌دست‌آمده تخمین داخل محدوده جدول است و تا یک رقم اعشار
گرد شده است. داده‌های تأخیر و زمان گذار به‌تنهایی برای بازسازی
دقیق شکل موج روی شبکه پیچیده \lr{RC} یا تحلیل دقیق نویز
اطلاعات محدودی دارند \cite[بخش \secnum{۴.۶.۲}، صفحه ۱۷۴]{weste}.

\subsection{مدل منبع جریان و شکل موج خروجی}
\label{sec:delay-current-source}

وقتی اتصال‌های بین گیت‌ها را با یک خازن معادل نمی‌توان به
دقت لازم توصیف کرد، شکل موج جریان راه‌انداز برای محاسبه
ولتاژ گره‌ها مفید است. در صورت‌بندی ایدئال یک
\lr{Current Source Model}، جریان خروجی به ولتاژ ورودی و
خروجی وابسته است. تحلیل‌گر جریان را همراه معادله‌های شبکه
\lr{RC} به کار می‌گیرد تا ولتاژ گره‌ها را در گذر زمان
به دست آورد \cite[بخش \secnum{۴.۶.۳}، صفحه ۱۷۴]{weste}.

برای دیدن سازوکار محاسبه، گره‌ای با ظرفیت ثابت \(C_L\) و
جریان خالص \(I_{\mathrm{net}}(t)\) به سمت گره را در نظر
بگیرید. در رابطه \eqref{eq:delay-step-kcl}، جریان با علامت
منفی باعث افت ولتاژ شد. در اینجا جهت مثبت را به سمت گره
می‌گیریم. جمع پیوسته جریان از زمان آغاز \(t_0\) تا \(t\)،
تغییر بار گره را می‌دهد. با تقسیم آن بر ظرفیت، ولتاژ خروجی
چنین به دست می‌آید.
\begin{equation}
  v_{\mathrm{out}}(t)=v_{\mathrm{out}}(t_0)
    +\frac{1}{C_L}\int_{t_0}^{t} I_{\mathrm{net}}(u)\,\mathrm{d}u.
  \label{eq:delay-current-source-step}
\end{equation}
در این رابطه \(u\) متغیر زمان در انتگرال است. انتگرال
جریان برحسب کولن و حاصل تقسیم آن بر ظرفیت برحسب ولت است.
اگر شبکه چند گره داشته باشد، جریان‌های مقاومت‌ها در تعادل بار
هر گره وارد می‌شوند. زمان عبور شکل موج از مرجع تعیین‌شده،
همراه زمان عبور ورودی، تأخیر انتشار را می‌دهد.

کتاب دو نمایش از داده‌های توصیف شکل موج را معرفی می‌کند.
\lr{Composite Current Source Model (CCSM)} جریان خروجی را
برحسب زمان، برای زمان گذار ورودی و بار معین، ذخیره می‌کند.
\lr{Effective Current Source Model (ECSM)} شکل موج ولتاژ
خروجی را ذخیره می‌کند. رابطه جریان و بار، مانند رابطه
\eqref{eq:delay-current-source-step} در حالت تک‌خازنی، ارتباط
این دو نوع داده را روشن می‌کند
\cite[بخش \secnum{۴.۶.۳}، صفحه ۱۷۴]{weste}.

\subsection{خطاهای رایج در مقایسه تأخیر}
\label{sec:delay-timing-pitfalls}

\begin{enumerate}
  \item \textbf{مقایسه عددهای تأخیر با بارهای متفاوت.}
    تأخیر یک مرحله \lr{Ring Oscillator} در مثال فصل \(2\tau\)
    و تأخیر \lr{FO4} برابر \(5\tau\) بود. اولی مربوط به
    بار یک ورودی هم‌اندازه و دومی مربوط به چهار ورودی هم‌اندازه
    است. برای مقایسه، بار و مرجع اندازه‌گیری را مشخص کنید
    \cite[بخش \secnum{۴.۷}، صفحات ۱۷۴ و ۱۷۵]{weste}.
  \item \textbf{بزرگ‌کردن گیت بدون محاسبه دوباره تأخیر مسیر.}
    بزرگ‌کردن یک گیت می‌تواند تأخیر خود آن گیت را کم کند و
    هم‌زمان بار مرحله پیشین را زیاد کند. در مرز دو بلوک،
    زمان رسیدن و قدرت راه‌اندازِ بلوک قبلی هر دو در
    بودجه زمان‌بندی اهمیت دارند
    \cite[بخش \secnum{۴.۷}، صفحه ۱۷۵]{weste}.
  \item \textbf{مقایسه مسیرها فقط با شمارش مراحل.}
    در \hyperref[sec:delay-path-stage-count]{مثال بار ۶۴ واحدی}،
    تقسیم بار بین چند مرحله، تأخیر را با وجود افزایش تعداد
    گیت‌ها کاهش داد. برای مقایسه مسیرها، هم \(F\) و هم \(P\)
    هر گزینه را محاسبه کنید
    \cite[بخش \secnum{۴.۷}، صفحه ۱۷۵]{weste}.
\end{enumerate}

\subsection{معیار \lr{FO4} و انتخاب مدل}
\label{sec:delay-model-selection}

در مرور تاریخی کتاب، زمان چرخه پردازنده‌ها برحسب تعداد
تأخیرهای \lr{FO4} گزارش شده است. برای نمونه، دوره کلاک
\lr{Pentium 4} در این مقایسه به حدود ۱۰ تأخیر \lr{FO4}
رسید. کتاب این کاهش را همراه با عمیق‌ترشدن \lr{Pipeline} و
هزینه توان بررسی می‌کند. اعداد شکل \secnum{۴.۳۸} برآوردهای
تاریخی‌اند و خود کتاب عدم قطعیت ناشی از تخمین \lr{FO4}
را یادآور می‌شود. بنابراین در مقایسه طرح‌ها، زمان چرخه را
همراه با معماری، بار و شرایط کاری تفسیر می‌کنیم
\cite[بخش \secnum{۴.۸}، شکل \secnum{۴.۳۸}،
صفحات ۱۷۵ و ۱۷۶]{weste}.

اگر مسئله تعیین زمان عبور یک سیگنال باشد، نخست
\hyperref[sec:delay-waveform]{مرجع اندازه‌گیری روی شکل موج}
را مشخص کنید. برای دیدن اثر مقاومت و ظرفیت گره‌ها، از
\hyperref[sec:delay-elmore]{مدل \lr{RC} و \lr{Elmore Delay}}
استفاده کنید. برای مقایسه سریع توپولوژی و اندازه‌گذاری،
\hyperref[sec:delay-path-model]{\lr{Logical Effort} مسیر}
اثر نوع گیت، بار و انشعاب را آشکار می‌کند. هرگاه زمان گذار
و شکل موج خروجی در نتیجه اثر تعیین‌کننده داشته باشند،
مدل زمانی کتابخانه و شبکه اتصال‌ها را در تحلیل وارد کنید
\cite[جمع‌بندی فصل \secnum{۴}، صفحه ۱۷۶]{weste}.
